% =============================================================================
\section{Results and Discussion}
% =============================================================================

% -----------------------------------------------------------------------------
\subsection{RQ1 -- Task-Level Capability Gap}
% -----------------------------------------------------------------------------

[To be completed after data collection.
Table~\ref{tab:metrics} will report mean ROUGE-1/2/L, F1, and rule-based pass
rates by task and condition with Wilcoxon signed-rank $p$-values and effect
sizes ($r = Z/\sqrt{N}$), Bonferroni-corrected across three tasks.
Expected pattern: LLM outperforms SLM on ROUGE for open-ended writing, where
output space is large and reference coverage is low; the gap narrows for
summarization (well-defined reference) and further for planning (rule-based
checklist scoring that rewards structure over fluency).]

\begin{table}[t]
  \caption{Objective output quality by task and condition (mean $\pm$ SD).
           \textit{[To be replaced with study data.]}}
  \label{tab:metrics}
  \small
  \begin{tabular}{llcc}
    \toprule
    Task & Metric & SLM & LLM \\
    \midrule
    Writing    & ROUGE-L   & -- & -- \\
    Summarize  & ROUGE-L   & -- & -- \\
    Planning   & Pass rate & -- & -- \\
    \bottomrule
  \end{tabular}
\end{table}

% -----------------------------------------------------------------------------
\subsection{RQ2 -- ``Good Enough'' Boundary}
% -----------------------------------------------------------------------------

[To be completed.
``Good enough'' is operationalized as: SLM metric $\geq 90\%$ of LLM metric
AND participant preference not significantly favoring LLM ($p > .05$,
binomial test).
Hypothesis: summarization and planning will meet the threshold; open-ended
creative writing will not, reflecting the higher output variance of generation
tasks without a reference ceiling.
If confirmed, this boundary becomes an actionable decision rule: for
structured everyday tasks, a privacy-preserving local SLM is a justified
default, while fluency-sensitive generation warrants the higher-resource LLM.]

% -----------------------------------------------------------------------------
\subsection{RQ3 -- Prompting-Skill Confound}
% -----------------------------------------------------------------------------

[To be completed.
Prompt logs will be coded for constraint completeness, specificity, and format
adherence (two coders, Cohen's $\kappa$ target $> 0.7$).
If the Goal Wizard and Prompt Plan Editor reduce within-condition variance in
constraint completeness---compared to a free-form baseline---this confirms that
scaffolding acts as a \textit{methodological control}, not only as a user aid.
The implication is that uncontrolled SLM--LLM comparisons in the literature
likely overstate the capability gap by attributing prompting-skill variance to
the model.]

% -----------------------------------------------------------------------------
\subsection{RQ4 -- Interface Affordances}
% -----------------------------------------------------------------------------

[To be completed.
NASA-TLX scores will quantify the workload cost of scaffolding.
Format-violation rates (checklist items missing, JSON malformed, constraint
words absent) will test whether step decomposition reduces constraint errors.
Expected: SLM+scaffold shows higher task demand but fewer format violations
than SLM without scaffold; a non-significant workload difference between
LLM (no scaffold) and SLM+scaffold would support the argument that
scaffolding compensates for model brittleness at acceptable cognitive cost.]

% -----------------------------------------------------------------------------
\subsection{Discussion}
% -----------------------------------------------------------------------------

\subsubsection{Everyday task needs are largely model-agnostic.}
The 65\% overlap between SLM and LLM task discourse confirms that users bring
the same goals to both model classes; what varies is the capability to satisfy
them and the friction of doing so.
This reframes the research question from ``what do people use SLMs for?''
to ``which tasks can SLMs serve well enough?''---a more tractable and
actionable question.

\subsubsection{Scaffolding as research infrastructure.}
The Goal Wizard and Prompt Plan Editor are designed to serve dual roles: a
user aid that reduces the Gulf of Envisioning for novices, and a
\textit{variance-reduction mechanism} that makes SLM--LLM comparisons
interpretable.
Making this dual role explicit is a methodological contribution: future
comparative studies can adopt or adapt Easy\_SLM's scaffolding layer as
a controlled prompting substrate, and the interface is open-sourced for reuse.

\subsubsection{Sustainable trade-offs made legible.}
The ``good enough'' boundary operationalized in RQ2 translates a technical
model-selection decision into user-facing information.
Displaying task-level sufficiency evidence in the interface---rather than
hiding it in deployment infrastructure choices---aligns with Crawford's call
to rematerialize AI trade-offs~\cite{crawford2024rematerializing}.
Users in the study who understand that SLM is ``good enough'' for their
planning task but not for their creative writing task are better equipped to
make intentional, sustainable choices about cloud dependency.

\subsubsection{Limitations.}
The formative Reddit study over-represents technically engaged users;
casual SLM adopters may report different task patterns.
ROUGE and rule-based pass rates are imperfect proxies for utility on
open-ended tasks and should be interpreted alongside preference and confidence
data.
Results are specific to the evaluated SLM [\textit{model name}, quantization];
different models may shift the capability gap substantially.
The within-subjects design introduces potential demand characteristics:
participants may moderate their behavior once they understand the comparison
purpose.

\subsubsection{Conclusion.}
Easy\_SLM pairs a formative task-discovery pipeline with a scaffolded
comparison interface to produce empirically grounded, confound-controlled
evidence about when SLMs are sufficient for everyday use.
The 19 shared tasks identified across SLM and LLM communities provide a
reusable task taxonomy for future work; the four design goals provide a
template for fair capability comparison interfaces; and the ``good enough''
framework offers a decision procedure for sustainable model selection that
is actionable for everyday users rather than exclusively for researchers
or system administrators.
